% Generated by scripts/assemble.py; corrections belong in source/corrections.json.
\examhead{149}{湖南师范大学 2026 年数学分析试卷}{141}
% p142-b22; source PDF page 142
\par\Needspace{5\baselineskip}
\blockanchor{p142-b22}
\majorquestion{一}填空题.\par
% p142-b23; source PDF page 142
\blockanchor{p142-b23}
\question{1}极限  $\lim_{x\to0}\frac{1}{x^{2}}\ln\frac{\sin x}{x}=$  \_\_\_\_.\par
% p142-b24; source PDF page 142
\blockanchor{p142-b24}
\question{2}极限  $\lim_{n\to\infty}\left(1-\frac{1}{2}+\cdots+(-1)^{n-1}\frac{1}{n}\right)=$  \_\_\_\_.\par
% p143-b00; source PDF page 143
\blockanchor{p143-b00}
\question{3}设 $f(x) = \frac{x^4}{1 - x^2}$，则 $f^{(2026)}(0) =$ \_\_\_\_.\par
% p143-b01; source PDF page 143
\par\Needspace{6\baselineskip}
\blockanchor{p143-b01}
\question{4}设 $f(x)$ 在 $(- \infty, + \infty)$ 上连续, $f'(0)$ 存在, 且\par
\nopagebreak[4]
% p143-b02; source PDF page 143
\blockanchor{p143-b02}
\displayformula{f (x) = x ^ {3} + x ^ {2} \lim _ {x \rightarrow 0} \frac {f (x)}{x} - x \int_ {0} ^ {1} f (x) \mathrm{d} x.}
% p143-b03; source PDF page 143
\blockanchor{p143-b03}
则  $f(x)=$  \_\_\_\_.\par
% p143-b04; source PDF page 143
\blockanchor{p143-b04}
\question{5}曲面 $x = \sqrt{2y^2 + 3z^2 - 4}$ 上垂直于直线 $\frac{x - 2}{1} = \frac{y - 3}{2} = \frac{z}{3}$ 的切平面方程为 \_\_\_\_.\par
% p143-b05; source PDF page 143
\blockanchor{p143-b05}
\question{6}设 $f(x)$ 二次可微, 且 $\sup_{x \in [0,1]} f''(x) \leq 2$, 若 $f(0) = f(1) = 2$, 且曲线 $y = f(x) - x^2$ 与直线 $y = 2 - x$ 交于点 $(x_0, y_0) (0 < x_0 < 1)$, 则对任意的 $x \in (0,1)$, $f(x) =$ \_\_\_\_.\par
% p143-b06; source PDF page 143
\blockanchor{p143-b06}
\question{7}设 p > 0, 广义积分  $\int_{0}^{+\infty}\frac{\sin x}{x^{p} + \sin x}dx$  条件收敛, 则实数 p 满足的条件是 \_\_\_\_.\par
% p143-b07; source PDF page 143
\blockanchor{p143-b07}
\question{8}级数 $\sum_{n=0}^{\infty} \frac{(-1)^n}{(2n)!} = \underline{\quad}$.\par
% p143-b08; source PDF page 143
\blockanchor{p143-b08}
\question{9}设曲线 $L: x^2 + y^2 = 16$, 取逆时针方向, 则 $\oint_{L} \frac{y \mathrm{~d}x - x \mathrm{~d}y}{x^2 + xy + y^2} =$\par
% p143-b09; source PDF page 143
\par\Needspace{6\baselineskip}
\blockanchor{p143-b09}
\question{10}计算曲面积分\par
\nopagebreak[4]
% p143-b10; source PDF page 143
\blockanchor{p143-b10}
\displayformula{I = \iint_ {S} \frac {2 x \mathrm{d} y \mathrm{d} z + (x + 2) ^ {2} \mathrm{d} x \mathrm{d} y}{\sqrt {x ^ {2} + y ^ {2} + z ^ {2}}}}
% p143-b11; source PDF page 143
\blockanchor{p143-b11}
其中 $S: z = -\sqrt{4 - x^2 - y^2}$, 取上侧, 则 $I =$ \_\_\_\_.\par
% p143-b12; source PDF page 143
\par\Needspace{5\baselineskip}
\blockanchor{p143-b12}
\majorquestion{二}计算题.\par
% p143-b13; source PDF page 143
\blockanchor{p143-b13}
\question{1}设 $f(x)$ 在 $[a, b]$ 上有一阶连续导数, 定义 $s_n = \int_0^1 f(x) \, \mathrm{d}x - \frac{1}{n} \sum_{k=1}^{n} f\left(\frac{k}{n}\right)$, 求 $\lim_{n \to \infty} n s_n$.\par
% p143-b14; source PDF page 143
\par\Needspace{6\baselineskip}
\blockanchor{p143-b14}
\question{2}设 $f(x, y)$ 在闭区域 $D = \{(x, y) \mid x^2 + y^2 \leq 1\}$ 上有二阶连续偏导数，且\par
\nopagebreak[4]
% p143-b15; source PDF page 143
\blockanchor{p143-b15}
\displayformula{\frac {\partial^ {2} f}{\partial x ^ {2}} + \frac {\partial^ {2} f}{\partial y ^ {2}} = \sin (x ^ {2} + y ^ {2}).}
% p143-b16; source PDF page 143
\blockanchor{p143-b16}
求 $\iint_{D}\left(x\frac{\partial f}{\partial x} +y\frac{\partial f}{\partial y}\right)\mathrm{d}x\mathrm{d}y.$\par
% p143-b17; source PDF page 143
\blockanchor{p143-b17}
\question{3}设函数 $I(\alpha) = \int_0^1\frac{\ln(1 + \alpha x)}{1 + x^2}\mathrm{d}x$ ($\alpha \geq 0$), 求 $I(1)$ 的值.\par
% p143-b18; source PDF page 143
\par\Needspace{5\baselineskip}
\blockanchor{p143-b18}
\majorquestion{三}证明题.\par
% p143-b19; source PDF page 143
\par\Needspace{6\baselineskip}
\blockanchor{p143-b19}
\question{1}设函数 $f(x), g(x)$ 在 $[a, b]$ 上连续，在 $(a, b)$ 内可导，且 $g'(x) \neq 0$，证明：存在 $\xi \in (a, b)$，使得\par
\nopagebreak[4]
% p143-b20; source PDF page 143
\blockanchor{p143-b20}
\displayformula{\frac {f ^ {\prime} (\xi)}{g ^ {\prime} (\xi)} = \frac {f (\xi) - f (a)}{g (b) - g (\xi)}.}
% p143-b21; source PDF page 143
\blockanchor{p143-b21}
\question{2}用有限覆盖定理证明：闭区间 $[a, b]$ 上的连续函数 $f(x)$ 一定在 $[a, b]$ 上一致连续.\par
% p143-b22; source PDF page 143
\par\Needspace{6\baselineskip}
\blockanchor{p143-b22}
\question{3}设函数 $f(x), g(x)$ 在 $[a, b]$ 上连续，且有相同的单调性，证明：\par
\nopagebreak[4]
% p143-b23; source PDF page 143
\blockanchor{p143-b23}
\displayformula{(b - a) \int_ {a} ^ {b} f (x) g (x) \mathrm{d} x \geq \int_ {a} ^ {b} f (x) \mathrm{d} x \int_ {a} ^ {b} g (x) \mathrm{d} x.}
% p143-b24; source PDF page 143
\blockanchor{p143-b24}
\question{4}设 $f(x)$ 定义在 $[-1, 1]$ 上, 且 $f''(0)$ 存在, 证明: $\sum_{n=1}^{\infty} f\left(\frac{1}{n}\right)$ 绝对收敛当且仅当 $f(0) = f'(0) = 0$.\par
